\documentclass[10pt,a4paper]{amsart}
\usepackage[margin=19mm]{geometry}
\usepackage{amsmath,amssymb,mathtools}
\usepackage[scr=rsfs,cal=euler]{mathalfa}
\usepackage{xfrac}
\usepackage[unicode,hidelinks,bookmarks=false]{hyperref}
\newcommand{\ie}{i.e.\ }
\newcommand{\cf}{cf.\ }
\newcommand{\resp}{respectively\ }
\newcommand{\Subsection}{\subsection}
\providecommand{\nobreakdash}{\nobreak\hbox{-}}
\setlength{\emergencystretch}{2em}
\multlinegap0pt
\usepackage{mathtools}
\usepackage{amssymb}
\usepackage{enumerate}
\usepackage{float}
\usepackage{xcolor}
\usepackage{tikz}
\usepackage{natbib}
\usepackage[normalem]{ulem}
\newtheorem{lemma}{Lemma}
\title{Spiking and Resetting}
\author{C\'edric Bernardin, Vsevolod Vladimirovich Tarsamaev}
\date{Source: arXiv:2512.19517v1 (CC BY 4.0)}
\begin{document}
\maketitle

\noindent\emph{Source and licence.} Spiking and Resetting. Version arXiv:2512.19517v1 (CC BY 4.0), distributed by arXiv under the Creative Commons Attribution 4.0 International licence, http://creativecommons.org/licenses/by/4.0/, as recorded in the arXiv licence metadata for this exact version and shown on the abstract page https://arxiv.org/abs/2512.19517v1. Full licence notice: \texttt{licenses/arxiv-2512.19517-CC-BY-4.0.txt}.\par\smallskip

\noindent\textit{Editorial scope.} Counted unit: the article's lemma lem:Vyepsilonstar (source0.tex lines 2396--2405). The article's complete original proof of it is delivered with it (source0.tex lines 2407--2464, one \texttt{proof} environment of 58 lines). No mathematics is written, repaired or re-derived: the statement and the proof are the article's own lines, in the article's own notation, and no retained line is substituted or re-flowed.  Retained uncounted as the source-local context the proof needs: lines 2244--2260.  Nothing was omitted from any retained span, so the package carries no omission record.  No retained line contains a citation, so the wrapper carries no bibliography and no bibliography entry.  No retained line contains a figure environment, an image-inclusion command or drawing code, so no image is delivered and the package declares no asset dependency.  Editorial formatting only, in addition: an AMS \texttt{amsart} preamble that restates the article's own declarations and macros used by the retained lines, namely the environments \texttt{lemma} on separate counters, together with the standard packages those lines need.  The retained environments print in this document as Lemma 1 in order of appearance, whereas the article prints the same items with the numbering of its own full document.  The complete native source of the article as distributed in the arXiv e-print for this exact version is bundled unchanged as \texttt{sources/191485/source0.tex}.
\begin{lemma}
\label{lem:expV}
There exists a constant $C>0$ such that for any $\alpha \in (0,1/2)$ and $x\in (0,1)$, the following holds
\begin{equation}
\label{eq:exp-Vvarepsilon}
\begin{split}
V^\varepsilon(x) &= - \log \varepsilon +  \log \Big(\frac{h(0)}{f(0)} \Big)  +  \log x + R_{V^\varepsilon} (x) 
\end{split}
\end{equation}
where 
\begin{equation}
%\label{eq:}
\begin{split}
\vert R_{V^\varepsilon} (x) \vert \le C \big[ \varepsilon^{1/2 - \alpha}/{\mathfrak  h}(x) + \varepsilon^{2\alpha} {\mathfrak h}(x)  \big] \ .
\end{split}
\end{equation}
\end{lemma}

\begin{lemma}
\label{lem:Vyepsilonstar}
We have that 
\begin{equation}
%
\begin{split}
\lim_{\varepsilon \to 0} e^{-V^\varepsilon (y_*^\varepsilon)} = 0 \ .
\end{split}
\end{equation}
\end{lemma}

\begin{proof}
Recall that \(y_*^\varepsilon = 1 - \varepsilon^\beta\) with \(\beta \in (0,1/2)\).  From Lemma \ref{lem:expV}, for any \(x \in (0,1)\) and \(\alpha \in (0,1/2)\),
\[
V^\varepsilon(x) = -\log \varepsilon + \log\left(\frac{h(0)}{f(0)}\right) + \log x + R_{V^\varepsilon}(x) \ ,
\]
with remainder satisfying
\[
|R_{V^\varepsilon}(x)| \leq C\bigl[\varepsilon^{1/2-\alpha}/\mathfrak{h}(x) + \varepsilon^{2\alpha}\mathfrak{h}(x)\bigr] \ ,
\]
where \(\mathfrak{h}(x) = \inf_{y\in[0,x]} h(y)\).

\bigskip

Take \(x = y_*^\varepsilon = 1 - \varepsilon^\beta\).   Since \(h\) is continuously differentiable, \(h(1)=0\) and $h'(1)<0$, there exists \(c>0\) such that for sufficiently small \(\varepsilon\),
\[
c \varepsilon^\beta \le \mathfrak{h}(y_*^\varepsilon) \le c^{-1} \varepsilon^\beta \ .
\]
Hence, there exists a constant $C>0$ such that
\[
|R_{V^\varepsilon}(y_*^\varepsilon)| \leq C\bigl[\varepsilon^{1/2-\alpha - \beta} + \varepsilon^{2\alpha+\beta}\bigr] \ .
\]
Choose \(\alpha = \beta/2\) (note that \(\beta<1/2\) implies \(\alpha\in(0,1/4)\)). Then
\[
\varepsilon^{1/2-\alpha-\beta} = \varepsilon^{1/2 - 3\beta/2}, \qquad \varepsilon^{2\alpha+\beta} = \varepsilon^{2\beta} \ .
\]
For \(\beta < 1/3\) (which is compatible with \(\beta\in(0,1/2)\)), both exponents are positive, so that 

\[
\lim_{\varepsilon \to 0} R_{V^\varepsilon}(y_*^\varepsilon) =0 \ .
\]

\bigskip

Now,

\[
V^\varepsilon(y_*^\varepsilon) = -\log \varepsilon + \log\left(\frac{h(0)}{f(0)}\right) + \log(1-\varepsilon^\beta) + o(1).
\]

Since \(\log(1-\varepsilon^\beta) = O(1)\), we obtain

\[
V^\varepsilon(y_*^\varepsilon) = -\log \varepsilon + O(1) \xrightarrow[\varepsilon\to0]{} +\infty.
\]

\bigskip

Consequently,

\[
\exp\bigl(-V^\varepsilon(y_*^\varepsilon)\bigr) = \exp\bigl(\log\varepsilon + O(1)\bigr) = \varepsilon \cdot e^{O(1)} \xrightarrow[\varepsilon\to0]{} 0.
\]

\bigskip

Thus, \(\exp(-V^\varepsilon(y_*^\varepsilon)) \to 0\) as required.

\end{proof}

\end{document}
